% !TEX root = ../main.tex
%%%%%%%%%%%%%%%%%%%%%%%%%%%%%%%%%%%%%%%%%%%%%%%%%%%%%%%%%%%%%%%%%%%%%%%
% Introduction
%%%%%%%%%%%%%%%%%%%%%%%%%%%%%%%%%%%%%%%%%%%%%%%%%%%%%%%%%%%%%%%%%%%%%%%

\section{Introduction}
\label{sec:intro}

The disaggregated Keynesian cross of \citet{Teglio2025} shows how heterogeneity in the marginal propensity to consume opens an output gap: when income accrues to agents who do not spend it, demand leaks out of the circular flow and the economy settles below capacity. The mechanism is a demand mechanism, and the supply side plays no active role: capacity is given.
That paper's conclusions name endogenous investment out of capitalists' accumulated wealth, affecting both the demand and the supply sides, as future work for the model, and venture a conjecture about its outcome; the present extension is that development, and \cref{sec:shape} returns to the conjecture.

This paper asks what happens when the supply side is made endogenous in the most direct way available, by letting firms invest out of their own profits and accumulate capital. The extension is not merely additive. Once capital accumulates and labor is hired endogenously, two channels that point in opposite directions become active simultaneously:

\begin{itemize}
  \item Supply: Higher retention $\rho$ means more investment, more capital, and higher productive capacity. Taken alone, this raises output.
  \item Demand: In a demand-constrained regime, more capital means fewer workers are needed to serve the same demand. Employment falls, the wage bill falls, and because workers have a higher propensity to consume than capitalists, demand falls with it.
\end{itemize}

Which channel dominates is the classical wage-led versus profit-led question \citep{BhaduriMarglin1990,LavoieStockhammer2013}. Posed as a question about a sign, it invites an answer in the form of a label. This paper's organizing finding is that the label is the wrong object: what the model determines is the shape of the response of output to retention, and the sign is a local reading of that shape.

\paragraph{Three contributions.}

First, $Y(\rho)$ is U-shaped, and the turn is where the economics is (\cref{sec:shape}). The response has a turning point $\rhostar$ that lies inside the swept support and rises with $\sigma$, from $0.42$ to $0.64$ across the range examined. The anchored retention rate, derived from the investment-rate anchor in \cref{sec:anchoredrho}, not assumed, lies to the left of that turn at every elasticity tested: the marginal effect of retention at the point where the data put the economy is negative, while its effect over a wide enough range is positive. Both are properties of one curve.
Summarized through a single slope, this structure presents as a sign frontier $\sigstar \approx 0.65$, a reduction we retain, and declare as one.

Second, the region in which retention depresses output is conditional on fiscal institutions (\cref{sec:government}). A balanced-budget unemployment benefit, at a replacement rate inside the OECD band, eliminates that region entirely, making retention expansionary at every elasticity tested. The leakage of the unemployed from the circular flow is the engine of the effect, and closing the leak removes it; a sign reported without its institutional setting is therefore incomplete.

Third, the wage curve is a source of fragility immune to demand instruments (\cref{sec:stress,sec:discussion}). Making wages flexible destabilizes the economy through a wage$\,\to\,U\,\to\,$capital-erosion feedback, and three separate stabilization hypotheses aimed at it, slower demand expectations, the unemployment benefit, and that benefit as a cushion against firm heterogeneity, are each falsified by the data designed to test them, all three through that same channel. The resistance to demand instruments is diagnostic rather than incidental: the feedback is a relative-price phenomenon, not a demand one. A price probe shows this directly: holding the numéraire fixed is itself a convention, and relaxing it to complete instantaneous pass-through switches the feedback off through a Pigou effect on real balances, so the finding is bracketed by the pass-through convention, with neither end anchored (\cref{sec:frontierlocal}).

A fourth contribution is methodological and kept deliberately off the main line of argument: the sign depends on the estimator used to measure it, a two-point chord and a whole-support slope disagreeing on $108$ identified design points under a verdict rule fixed in advance (\cref{sec:estimator}). The same analysis shows that the frontier does not survive marginalization, viability is governed almost entirely by the depreciation rate rather than by $\sigma$, and at the rate the data imply, no sampled configuration remains viable (\cref{sec:sa}).

\paragraph{A note on what this paper reports.}
An unusual proportion of what follows is negative: three hypotheses are falsified, one headline does not generalize, and one measurement instrument turns out to carry part of its own answer. This is by design, under a declared rule on what may be reported as a result (\cref{sec:method}); the claims withdrawn under that rule are listed in \cref{app:retractions} rather than quietly dropped.

The remainder is organized as follows.
\Cref{sec:literature} places the model in the literature and \cref{sec:model} presents it. \Cref{sec:calibration} documents parameter anchoring and \cref{sec:method} the experimental protocol.
\Cref{sec:shape} to \cref{sec:sa} report the four blocks of results, \cref{sec:discussion} draws the cross-cutting readings, \cref{sec:limits}
states the limitations, and \cref{sec:conclusion} concludes.
\Cref{app:validation} collects the validation record.
